\section{Literature Review}
\label{sec:literature-review}
The relevant literature connects stream-processing foundations, numerical resource optimization, stateful reconfiguration, fault recovery, performance evaluation, cost-aware information acquisition, and agentic operational diagnosis. These strands already provide many of the mechanisms that an adaptive streaming controller would require. The purpose of this review is to distinguish those inherited capabilities from the proposed evidence-acquisition question, rather than to treat their combination as automatically novel. The bibliography contains forty-one published papers. Table~\ref{tab:related-work} positions the ten closest technical precedents; the remaining foundational and cross-cutting studies are integrated thematically below. Comparisons concern the inspected mechanisms and evaluation scope, not controlled reruns of the published systems.

\begin{table*}[t]
\caption{Closest technical precedents for the planned EcoStream study. The final column records the comparison obligation; it does not assert that a cited system lacks all other capabilities.}
\label{tab:related-work}
\centering
\footnotesize
\setlength{\tabcolsep}{4pt}
\renewcommand{\arraystretch}{1.13}
\begin{tabularx}{\textwidth}{@{}p{0.13\textwidth} p{0.21\textwidth} p{0.25\textwidth} >{\raggedright\arraybackslash}X@{}}
\toprule
\textbf{Work} & \textbf{Primary decision} & \textbf{Evidence strategy} & \textbf{Comparison obligation} \\
\midrule
Phoebe~\cite{geldenhuys2022phoebe} & Resource configuration under latency and recovery constraints & Profiling, workload forecasting, and analytical models & Recovery-aware numerical baseline \\
Demeter~\cite{geldenhuys2024demeter} & Multi-configuration resource optimization & Forecasting plus selective profiling when model knowledge is insufficient & Numerical and information-acquisition comparator \\
Scarf~\cite{liu2026scarf} & Multi-objective Flink configuration tuning & Runtime state and transferable learned models & Learning-based tuning comparator \\
STREAMLINE~\cite{pedratscher2025streamline} & Ensemble-wide tuning, scheduling, and updates & Workload prediction and evolutionary search & Holistic adaptation precedent \\
CAOM~\cite{tan2025caom} & Operator migration timing & Runtime bottleneck metrics and interruption cost & Transition- and backlog-aware decision precedent \\
DRRS~\cite{qing2025drrs} & Execution of fine-grained stateful scaling & Routing, ordering, and state-migration information & Transition-cost evidence, not a diagnosis baseline \\
RCAgent~\cite{wang2024rcagent} & Root-cause diagnosis for cloud and Flink incidents & Agent-selected diagnostic tools and explicit finalization & Conventional agentic acquisition and stopping comparator \\
ORACL~\cite{bai2026oracl} & Root-cause reasoning and resource recommendations & Aggregated telemetry, forecasting, and guarded execution & Closest diagnosis-to-allocation overlap \\
TAMO~\cite{zhang2025tamo} & Tool-assisted multimodal diagnosis & Metrics, logs, traces, and specialized tools & Multimodal agentic comparator; final method review pending \\
CausalRCA~\cite{xin2023causalrca} & Fine-grained root-cause ranking & Learned causal graph over operational metrics & Non-LLM diagnostic comparator; action effects remain separate \\
\bottomrule
\end{tabularx}
\end{table*}

Several surveys provide complementary foundations. Liu and Buyya organize resource management around provisioning, operator parallelization, and scheduling under service and cost constraints~\cite{liu2020resource}. Fragkoulis et al. connect stream-processing evolution to state, ordering, fault tolerance, elasticity, and reconfiguration~\cite{fragkoulis2024evolution}, while Cardellini et al. classify runtime adaptation strategies and open challenges~\cite{cardellini2022adaptation}. Wang et al. review fault-tolerance mechanisms across stream-processing systems~\cite{wang2021faulttolerance}; Zhang et al. separately survey transactional stream processing and the consistency implications of combining transactions with continuous execution~\cite{zhang2023transactional}. Together, these sources motivate treating adaptation as an established systems problem constrained by semantics, state, recovery, scheduling, and service quality. As surveys, they inform terminology and coverage rather than constitute executable controller baselines.

\subsection{Numerical Autoscaling and Configuration Optimization}
Phoebe combines workload forecasting, parallel profiling, and quality-of-service models to tune streaming jobs under latency and recovery requirements~\cite{geldenhuys2022phoebe}. Its recovery discussion establishes an important constraint: capacity sufficient to match normal input can still be insufficient to clear the backlog after a failure. Daedalus instead builds processing-capacity models from runtime observations and uses forecasting for proactive horizontal scaling~\cite{pfister2024daedalus}. Its comparisons illustrate that reducing resources can involve a latency trade-off, rather than improving every service metric simultaneously.

Demeter expands the decision space to worker count, CPU, memory, processing slots, and checkpoint interval, combining forecasting with multi-objective Bayesian optimization~\cite{geldenhuys2024demeter}. It already selects additional configuration-profiling runs according to information gain and skips unnecessary profiling. Consequently, selective acquisition of information in general is not an unoccupied research area. The candidate distinction for EcoStream concerns operational diagnostic observations that separate competing bottleneck explanations, rather than additional executions intended to learn a configuration's performance.

More recent work broadens adaptation beyond isolated job parameters. STREAMLINE coordinates workload prediction, evolutionary configuration tuning, and operator scheduling across pipeline ensembles, with a state-preserving update mechanism described in its publisher abstract~\cite{pedratscher2025streamline}. MUDAP exposes service- and resource-level elasticity parameters, while its RASK agent uses structural knowledge, regression, and a constrained solver to coordinate competing services~\cite{sedlak2026mudap}. Some MUDAP actions alter service quality or model size. Such flexibility is valuable in its setting, but cannot be assumed equivalent to changing resources while preserving an exact streaming query. EcoStream must therefore compare controllers under the same permitted actions and semantics, including a strong low-overhead numerical alternative.

Scarf provides a particularly relevant recent Flink comparator. It combines workload-sensitivity-based knob selection, multi-objective reinforcement learning, and transferable graph-based models to tune throughput and resource usage~\cite{liu2026scarf}. Its publisher-deposited abstract describes offline learning of configuration trade-offs and online selection. Therefore, learning resource--performance trade-offs or transferring configuration knowledge is not the proposed novelty. Any adapted Scarf baseline must disclose differences in knobs, training effort, and available workloads.

Prediction and heterogeneous resource choice also have strong numerical precedents. MMP-SPS separates monitoring, predictor selection, and adaptation across multiple control windows, comparing prediction-error and bandit-based selection in Apache Storm~\cite{wladdimiro2026mmpsps}. Its evaluation also distinguishes resource savings from service degradation. Russo Russo et al. coordinate application- and operator-level controllers using Bayesian optimization and reinforcement learning for heterogeneous nodes~\cite{russorusso2023hierarchical}. Li et al.'s CEBFD approach combines resource-demand prediction with cost-aware executor placement for Spark Streaming~\cite{li2024cebfd}. The latter's micro-batch execution is not interchangeable with the proposed Flink setting; its relevance is the separation of demand prediction from placement cost.

ProOSStream combines offline scheduling optimization, federated reinforcement learning for online adaptation, and workload forecasting in edge--cloud environments~\cite{lajili2026proosstream}. It broadens the comparison beyond centralized scaling, but importing its distributed architecture into a single-host study would not establish equivalent deployment evidence. These works support testing whether diagnostic acquisition contributes beyond accurate prediction and competent numerical scheduling, with learning and controller overhead included.

Vogel et al. demonstrate online and transparent self-adaptation of parallel stream patterns, reinforcing that runtime adjustment without an agent is an established alternative~\cite{vogel2023selfadaptation}. Li et al. jointly consider scheduling cost and network topology in distributed stream computing~\cite{li2026costtopology}. These studies narrow the proposed contribution: EcoStream cannot claim novelty for adapting parallelism or considering topology. Its comparison must instead establish whether diagnostic evidence changes action selection beyond what a capable numerical controller can infer from its monitored state.

Evaluation methodology is also part of the comparison. ShuffleBench isolates large-scale data-shuffling behavior and compares distributed stream-processing frameworks under controlled workloads~\cite{henning2024shufflebench}. Although it is not an autoscaler, it illustrates why throughput or latency from one engine and topology should not be treated as a universal property. EcoStream experiments should report framework, serialization, partitioning, state, offered load, and network conditions, and should avoid attributing engine-specific effects to the controller.

\subsection{Partitioning, Scheduling, and Stateful Reconfiguration}
Performance deterioration can originate in data distribution, communication, or the mechanics of adaptation, not only insufficient CPU allocation. Ra-Stream combines resource scaling with communication-aware subgraph partitioning and fine-grained task scheduling~\cite{fan2026rastream}. SAP-DPO responds to fluctuating input with stream-aware partitioning, adaptive update cycles, and sliding-window parameter adjustment~\cite{ren2026sapdpo}. These approaches motivate separating capacity pressure from distribution-related problems when interpreting streaming telemetry. They also constrain a broad novelty claim about adaptive monitoring or update frequency.

CAOM detects bottleneck operators from runtime metrics and selects migration timing while considering interruption-related data accumulation~\cite{tan2025caom}. Thus, bottleneck-sensitive interventions and time-dependent migration costs already have direct streaming precedents. DRRS addresses the execution of stateful scaling through re-routing, record reordering, and subdivision of migrating state~\cite{qing2025drrs}. Its role differs from selecting a configuration: it reduces disruption when a scaling operation has already been chosen.

This distinction matters for EcoStream's transition model. Not every stateful adjustment should be assigned a mandatory full-restart cost; the model must reflect the actual execution path of the selected engine and configuration. Likewise, detecting skew or an impaired worker is not successful remediation unless the shared action catalogue contains an intervention that can address it. Cases outside that catalogue require an explicit no-change or escalation outcome, not an assumed successful repair.

Additional scheduling and recovery work makes this boundary more concrete. Ms-Stream coordinates grouping, partitioning, and placement to mitigate task- and node-level stragglers~\cite{wu2025msstream}. Ns-Stream instead prioritizes near-source processing through dynamic grouping, with its published implementation using Storm~\cite{wu2025nsstream}. Both inform the distinction between capacity and distribution bottlenecks; neither mechanism can be assumed available in an unmodified Flink executor. Incidents requiring unsupported redistribution should be excluded from a remediation claim or explicitly escalated.

Recovery scope is itself a design variable. Takdir et al. implement partial snapshots and local recovery in Flink to restore a subset of operators rather than the entire job~\cite{takdir2025localrecovery}. A-FP4S distributes state fragments over a peer-to-peer overlay and adapts parallel recovery to hardware heterogeneity~\cite{xu2023afp4s}. These mechanisms establish alternatives to global restoration. EcoStream would select among supported interventions, not claim to invent these recovery mechanisms; its measurements must use the actual restart, replay, and backlog costs observed in the deployed configuration.

Checkpoint protocol and execution architecture further affect those costs. CheckMate compares coordinated, uncoordinated, and communication-induced checkpointing for streaming dataflows under multiple workload conditions~\cite{siachamis2024checkmate}. Scabbard studies single-node fault-tolerant stream processing and demonstrates a recovery design whose assumptions differ from distributed Flink recovery~\cite{theodorakis2021scabbard}. These works caution against assigning one generic transition penalty to all interventions. EcoStream's cost model must be calibrated for the actual protocol and deployment used in each experiment, and its conclusions should not be generalized to recovery architectures that were not tested.

\subsection{Agentic Diagnosis and Tool-Oriented Reasoning}
RCAgent is a particularly close precedent because autonomous operational diagnosis has already been applied to Alibaba Cloud Flink jobs~\cite{wang2024rcagent}. Its controller gathers information through tools, manages lengthy observations using snapshot keys, repairs malformed outputs, and can finalize a diagnosis. The existence of autonomous tool selection, efficient context management, or a stopping action is therefore not sufficient to distinguish EcoStream. The proposed difference concerns the criterion for acquiring evidence and stopping, and its effect on executed interventions.

ORACL couples language-model root-cause reasoning with resource-allocation recommendations~\cite{bai2026oracl}. Its inspected manuscript delegates numerical forecasting to a separate component and mediates actions through a verification interface. Diagnosis before allocation and separation of numerical prediction from language interpretation are consequently inherited design choices. EcoStream's planned comparison must ask whether another observation changes the preferred feasible action enough to justify its cost, rather than whether an agent can generate an allocation recommendation.

Scientific tool agents provide supporting execution patterns, but not evidence of streaming-control performance. SciToolAgent uses knowledge-graph-assisted retrieval for scientific tool selection and execution, including safety checks~\cite{ding2025scitoolagent}. BioMedAgent separates planning, coding, and execution, using interactive exploration and retrieval of successful experiences to construct biomedical workflows~\cite{bu2026biomedagent}. Their relevance is bounded to tool contracts, feedback, and experience organization. A large tool graph or additional agent roles should not be introduced into a small operational catalogue without evidence that the extra complexity improves decisions. Transfer from biomedical or scientific workflow success to latency-sensitive streaming control remains a hypothesis.

More recent diagnosis work strengthens the need for a precise comparison. TAMO's published scope explicitly combines tool-assisted LLM diagnosis with multimodal observations in cloud-native systems~\cite{zhang2025tamo}. LocaleXpert converts metrics, logs, and traces into language-based descriptions and coordinates specialist agents to localize and interpret failures~\cite{zhong2026localexpert}. MA-RCA separates diagnostic roles and introduces retrieval grounding and runtime hypothesis validation~\cite{fu2026marca}. Consequently, tool use, multimodal evidence, specialist agents, and hypothesis checking are established capabilities, not sufficient contribution statements on their own. Detailed reconciliation with TAMO's inaccessible final method remains necessary.

CausalRCA provides a non-LLM alternative, learning a weighted causal graph from metrics and ranking fine-grained candidate causes~\cite{xin2023causalrca}. Such diagnostic ranking does not by itself measure the benefit of an executed repair, and an observational causal model still requires its assumptions to be checked. The proposed study should separately measure diagnosis quality, action selection, investigation overhead, and post-intervention service recovery. It should not interpret a more convincing explanation as evidence of a better operational outcome.

Multimodal diagnosis predates the newest LLM agents. Nezha performs interpretable fine-grained root-cause analysis over multimodal observability data~\cite{yu2023nezha}. DiagFusion combines metrics, logs, traces, deployment information, event representations, and graph learning to localize faulty instances and classify failure type~\cite{zhang2023diagfusion}. These systems support the premise that different modalities can provide complementary evidence, but they also provide strong non-agentic comparators. EcoStream must therefore test the value of sequential, action-sensitive acquisition against fixed multimodal fusion, not merely against a single-metric threshold.

Cost-aware sequential observation is also an established problem outside operational diagnosis. The active-feature-acquisition literature formalizes the choice of which observation to obtain next and when to stop, balancing predictive loss against acquisition cost through model-based, model-free, and hybrid policies~\cite{aronsson2026afa}. EcoStream therefore cannot claim novelty for cost-aware acquisition or stopping in general. Its narrower question is whether an operational observation is worth collecting because its possible outcomes can change the preferred feasible intervention after state-transition and recovery costs are included.

\subsection{Sustainability and Measurement Boundaries}
Energy-aware scheduling is also established in Flink. Li et al. propose an SLA-based scheduling strategy together with coordinated inter-node and intra-node load balancing~\cite{li2025energyflink}. The reviewed abstract supports its scheduling motivation, but does not provide sufficient instrumentation detail for an independent physical-energy assessment. Resource allocation and energy consumption must therefore remain distinct quantities in comparisons.

Colangelo et al. demonstrate grid-responsive orchestration of AI workloads, using power controls, pausing, and resource reallocation while respecting designated performance-flexibility tiers~\cite{colangelo2026grid}. This provides a physical-measurement precedent, not a direct streaming-controller baseline: some AI-job service commitments permit throughput reductions, whereas EcoStream's proposed exact-query evaluation preserves semantics and offered load. Allocated core-seconds, consumed CPU, instantaneous power, total energy, and carbon emissions answer different questions. Without suitable physical instrumentation, the present study will target resource efficiency and service quality; sustainability remains motivation rather than a measured energy-saving claim.

\subsection{Comparative Gap and Proposed Study Question}
The reviewed systems already cover proactive adaptation, constrained optimization, selective profiling, autonomous diagnosis, and recovery- or migration-aware execution. Their strengths motivate a specific unresolved question for this study: under what conditions does acquiring an additional operational observation improve a stateful intervention enough to outweigh collection and waiting costs? This is a research question, not an assertion that every prior controller lacks decision-sensitive acquisition.

EcoStream would investigate an acquisition and stopping policy based on expected intervention loss. Its intended evaluation must separate this policy from the contributions of language interpretation and transition-cost modelling, using matched diagnostic tools and actions. A conventional diagnostic agent and deterministic processing of the same evidence are necessary controls. The proposed benefit would be weakened if cheap telemetry already determines the best action, if deterministic rules match the agent, or if investigation delays erase resource savings. These possible outcomes are part of the study, not exceptions to be removed after testing. Priority for the specific mechanism remains provisional until the closest final versions and implementations are reconciled.

\subsection{Scope of the Current Literature Evidence}
This is a targeted comparison, not an exhaustive systematic review. The forty-one-paper bibliography emphasizes recent method papers while retaining older published work when it directly establishes a resource-management, fault-tolerance, self-adaptation, benchmarking, multimodal-diagnosis, or evidence-acquisition foundation. Some records were reviewed through publisher abstracts or author manuscripts; inaccessible method passages and differences from final published versions remain unresolved. TAMO is included using its published identity and title-level scope only, without importing uninspected preprint details. The bibliography identifies published versions, not ordinary preprints. Mechanism descriptions are limited to inspected support, and this draft neither reproduces literature-reported results nor presents them as EcoStream findings. The project's reference ledger records access and verification status; pending final-version and notice checks prevent a submission-readiness claim.
